% Lösungen Einheit 1
\einheitenkopf{Einheit 1 von 4 · Wurzelziehen und Lösbarkeit}
\erg{1}{a) quadratisch \quad b) linear \quad c) quadratisch \quad d) linear \quad e) quadratisch}
\erg{2}{a) positiv \quad b) negativ \quad c) null \quad d) positiv \quad e) negativ}
\erg{3}{a) $x_1 = 2$, $x_2 = -2$ \quad b) $x_1 = 9$, $x_2 = -9$ \quad c) $x_1 = 1$, $x_2 = -1$ \quad d) $x_1 = 10$, $x_2 = -10$ \quad e) $8\,$m (nur die positive Lösung ist eine Länge) \quad f) $x_1 = \sqrt{11} \approx 3{,}32$, $x_2 = -\sqrt{11} \approx -3{,}32$ \quad g) keine Lösung, ein Quadrat ist nie negativ}
\erg{4}{a) $x^2 = 25$; $x_1 = 5$, $x_2 = -5$ \quad b) $x^2 = 64$; $x_1 = 8$, $x_2 = -8$ \quad c) $3x^2 = 12$, $x^2 = 4$; $x_1 = 2$, $x_2 = -2$ \quad d) $-x^2 = -16$, $x^2 = 16$; $x_1 = 4$, $x_2 = -4$ \quad e) $r^2 = \frac{V}{\pi \cdot h} = \frac{1000}{12\pi} \approx 26{,}53$; $r \approx 5{,}15\,$cm (nur die positive Lösung ist ein Radius) \quad f) $x^2 = -4$: keine Lösung \quad g) $x^2 = 0$: genau eine Lösung, $x = 0$}
\erg{5}{a) $x_1 = 4$, $x_2 = -2$ \quad b) $x_1 = 6$, $x_2 = -2$ \quad c) $x_1 = 6$, $x_2 = 4$ \quad d) $x + 2 = 3$ oder $x + 2 = -3$; $x_1 = 1$, $x_2 = -5$ \quad e) $x - 4 = 0$; $x = 4$ (genau eine Lösung)}
\erg{6}{a) $4 = 4$ (wA) \quad b) $8 \neq 10$ (fA) \quad c) $32 = 32$ (wA) \quad d) $5 = 5$ (wA) \quad e) $18 \neq 16$ (fA) \quad f) $9 \neq 1$ (fA) \quad g) $16 = 16$ (wA)}
\erg{7}{a) zwei: $x_1 = 0$, $x_2 = 2$ \quad b) eine: $x = 1$ \quad c) keine (die Gerade $y = -3$ liegt unter dem Scheitel) \quad d) keine: $S({-2}|1)$ liegt über $y = 0$, Parabel nach oben geöffnet \quad e) zwei: $S(1|3)$ liegt über $y = -2$, Parabel nach unten geöffnet \quad f) eine: $S(3|{-5})$ liegt auf $y = -5$}
\erg{8}{a) jede negative Zahl, z.\,B. $x^2 = -2$ \quad b) $x^2 + 7 = 7$ \quad c) wahr, $x = 6$ ist die einzige Lösung; falsch, die Lösungen sind $-2$ und $-8$; z.\,B. $x^2 + 4 = 0$}
\erg{9}{a) Die negative Lösung fehlt: $x_1 = 4$, $x_2 = -4$ \quad b) $x^2 = 81$; $x_1 = 9$, $x_2 = -9$ \quad c) Vorzeichenfehler, $6$ muss abgezogen werden: $x_1 = 2 - 6 = -4$, $x_2 = -2 - 6 = -8$ \quad d) $x + 5 = 2$ oder $x + 5 = -2$; $x_1 = -3$, $x_2 = -7$}
\erg{10}{a) $\sqrt{30}$ und $-\sqrt{30}$ ergeben quadriert beide $30$. \quad b) Ein Quadrat ist nie negativ, also kann $x^2$ nicht $-30$ sein. \quad c) Nein: $\sqrt{0} = 0$, also $x - 2 = 0$ und $x = 2$ ist genau eine Lösung.}
